\section{Conceptual Soundness}\label{ConceptualSoundness}
\pagestyle{mypagestyle}
\begin{spacing}{1.75}
	
	\MUFGfont{
		\subsection{Objective and evaluation criteria}
		Conceptual soundness is assessed by checking whether:
		\begin{itemize}
			\item the FX dynamics are mathematically well-defined and consistent with the broader simulation framework;
			\item drift specification is consistent with arbitrage-free FX relations (forward parity under the chosen numeraire/measure);
			\item modelling simplifications are appropriate for PFE usage and horizons (including MPOR sensitivity);
			\item model limitations are clearly documented and mitigated by controls, overlays, or monitoring.
		\end{itemize}
	}
	
	\MUFGfont{
		\subsection{Arbitrage-free consistency of forward-implied drift}
		The mean function $g_k^{FX}(t)$ is constructed from discount factors, which is consistent with covered interest parity:
		\[
		F_{0,t}=S_0\frac{DF_{\text{FOR}}(0,t)}{DF_{\text{USD}}(0,t)}
		\]
		under the assumed quoting convention. The model uses $g_k^{FX}(t)$ as the deterministic component such that
		$\mathbb{E}[Y_k^{FX}(t)]=g_k^{FX}(t)$, ensuring the simulated spot has an expected path matching the initial forward curve.
		
		\textbf{Validation expectation:} implement numerical martingale/consistency checks:
		\begin{itemize}
			\item $\mathbb{E}\!\left[\frac{DF_{\text{USD}}(0,t)}{DF_{\text{FOR}}(0,t)}Y_k^{FX}(t)\right]\approx Y_k^{FX}(0)$ (appropriate discounted expectation);
			\item cross-rate identity: for currencies $a,b$ with USD legs,
			$Y_{a/b}(t)=Y_{a/USD}(t)/Y_{b/USD}(t)$ holds pathwise.
		\end{itemize}
	}
	
	\MUFGfont{
		\subsection{Distributional assumptions and tail behaviour}
		Under constant volatility GBM, log returns are normal with variance proportional to horizon.
		This implies:
		\begin{itemize}
			\item no volatility clustering (returns are independent conditional on $\sigma$);
			\item symmetric log-return distribution (no skew), and lognormal spot distribution;
			\item tails lighter than what is often observed in FX during stress.
		\end{itemize}
		These are common simplifications in exposure simulation engines but should be explicitly stated as limitations, with compensating controls
		(e.g., stress scaling, conservative overlays, or scenario add-ons).
	}
	
	\MUFGfont{
		\subsection{Horizon alignment: 5-day return calibration vs PFE drivers}
		The volatility calibration uses 5-day log returns, which may be intended to align with MPOR-like horizons.
		However, the PFE horizon and simulation time-step may differ, and the mapping should be explicit:
		if simulation uses daily steps, the implied daily vol is $\sigma_{\text{1d}}=\sigma_{\text{5d}}/\sqrt{5}$ under i.i.d. assumptions.
		Validation should include evidence that this mapping is stable across windows and regimes.
	}
	
\end{spacing}
\clearpage
